% syllabus.tex — Bilingual course syllabus for "Quantum Scientific Computing".
% Compile with XeLaTeX (ctexart auto-detects the macOS Chinese fontset).
%   xelatex syllabus.tex
% Content mirrors the SJTU graduate-course application form
% (forms/附件1：上海交通大学研究生课程开设申请表.docx).
\documentclass[UTF8,11pt]{ctexart}
\usepackage[margin=1in]{geometry}
\usepackage{booktabs}
\usepackage{tabularx}
\usepackage{array}
\usepackage{enumitem}
\usepackage{xcolor}
\usepackage{hyperref}
\hypersetup{colorlinks=true,linkcolor=blue!55!black,urlcolor=blue!55!black}

\newcolumntype{Y}{>{\raggedright\arraybackslash}X}
\newcommand{\hh}[1]{\noalign{\smallskip}\multicolumn{1}{@{}l@{}}{}&#1} % (unused helper placeholder)
\setlength{\extrarowheight}{2pt}

\title{\bfseries 量子科学计算 \\ Quantum Scientific Computing \\[0.4em]
  \large 课程教学大纲 \quad Course Syllabus}
\author{上海交通大学 数学科学学院 \\ School of Mathematical Sciences, Shanghai Jiao Tong University}
\date{2026--2027 学年 秋季学期 \quad Fall 2026--2027}

\begin{document}
\maketitle

% ============================================================== 课程基本信息
\section*{一、课程基本信息 \quad Course Information}
\addcontentsline{toc}{section}{课程基本信息 / Course Information}

\begin{tabularx}{\textwidth}{@{}l Y l Y@{}}
\toprule
\textbf{课程名称 Course} & 量子科学计算 / Quantum Scientific Computing
  & \textbf{学分 Credits} & 3 \\
\textbf{总学时 Hours} & 48（1 学分 $\ge$ 16 课时）
  & \textbf{开课季节 Semester} & 秋季学期 / Fall \\
\textbf{是否跨学期 Cross-sem.} & 否 / No
  & \textbf{课程性质 Nature} & 专业课 / Specialized Course \\
\textbf{课程分类 Type} & 全日制课程 / Full-time
  & \textbf{课程类别 Category} & 专业选修课 / Specialized Elective \\
\textbf{课程层次 Level} & 硕博共用 / All Graduates
  & \textbf{授课语言 Language} & 中英文并行 / Chinese \& English \\
\textbf{上课方式 Method} & 课堂教学 / In-class
  & \textbf{成绩记录 Grade} & 等级制 / Letter grading \\
\textbf{考试类型 Exam} & 论文 / Essay
  & \textbf{开课单位 School} & 数学科学学院 / Math. Sciences \\
\bottomrule
\end{tabularx}

\medskip
\noindent\textbf{负责教师 Instructor in charge:}\ 张镭 (Lei Zhang)\quad
工号/ID 10900\quad 数学科学学院与自然科学研究院 (School of Mathematical Sciences \&
Institute of Natural Sciences)\quad \href{mailto:lzhang2012@sjtu.edu.cn}{lzhang2012@sjtu.edu.cn}

% ==================================================================================
\section*{二、课程简介 \quad Course Description}
\addcontentsline{toc}{section}{课程简介 / Course Description}

\paragraph{中文}
本课程是一门面向科学计算应用的量子计算交叉课程，主要介绍量子线路、量子算法、量子线性代数、
量子模拟以及量子算法在偏微分方程、动力系统、金融、材料和机器学习中的应用。课程目标是帮助学生
掌握量子态、量子门、相位估计、QSP、QSVT、HHL、薛定谔化等基本思想与方法，并理解其在高维科学
计算中的潜在优势与局限。学生应具备线性代数、微积分、微分方程、数值方法和基础编程知识；量子力学
有助于理解部分内容，但并不是必须先修课程。

\paragraph{English}
This course is an interdisciplinary quantum computing course oriented toward scientific
computing applications. It introduces quantum circuits, quantum algorithms, quantum linear
algebra, quantum simulation, and applications of quantum algorithms to partial differential
equations, dynamical systems, finance, materials science, and machine learning. The goal is to
help students understand the basic ideas and methods of quantum states, quantum gates, phase
estimation, QSP, QSVT, HHL, and Schr\"odingerization, and to recognize their potential
advantages and limitations in high-dimensional scientific computing. Students are expected to
have background knowledge in linear algebra, calculus, differential equations, numerical
methods, and basic programming. Quantum mechanics is helpful for understanding some parts of
the course, but it is not a required prerequisite.

% ==================================================================================
\section*{三、教学目标 \quad Learning Outcomes}
\addcontentsline{toc}{section}{教学目标 / Learning Outcomes}
完成本课程后，学生应能够 \textit{(upon completion, students will be able to)}：
\begin{enumerate}[leftmargin=1.6em,itemsep=1pt]
  \item 使用量子态、量子门与量子线路模型描述量子计算，并分析其复杂度。\\
        \textit{Describe computation with quantum states, gates, and circuits, and analyze complexity.}
  \item 推导并运用核心量子算法原语：QFT、相位估计、Hadamard test、振幅估计。\\
        \textit{Derive and apply core primitives: QFT, phase estimation, Hadamard test, amplitude estimation.}
  \item 掌握量子数值线性代数工具链：block encoding、LCU、QSP/QSVT、HHL。\\
        \textit{Master the QNLA toolchain: block encoding, LCU, QSP/QSVT, and HHL.}
  \item 运用 Hamiltonian 模拟与薛定谔化方法求解微分方程与动力系统问题。\\
        \textit{Use Hamiltonian simulation and Schr\"odingerization to solve differential equations and dynamics.}
  \item 评估量子科学计算在金融、材料、机器学习等高维问题中的优势、局限与实现代价。\\
        \textit{Assess the advantages, limits, and costs of quantum scientific computing in high-dimensional applications.}
\end{enumerate}

\paragraph{先修要求 Prerequisites.}
线性代数、微积分、微分方程、数值方法、基础编程；量子力学非必需。\quad
\textit{Linear algebra, calculus, differential equations, numerical methods, basic programming;
quantum mechanics is not required.}

% ==================================================================================
\section*{四、教学大纲（中文） \quad Syllabus (Chinese)}
\addcontentsline{toc}{section}{教学大纲（中文）}
\renewcommand{\arraystretch}{1.2}
\begin{tabularx}{\textwidth}{@{}p{3.2cm} Y c p{3.1cm}@{}}
\toprule
\textbf{章节} & \textbf{主要内容} & \textbf{课时} & \textbf{教学方式} \\
\midrule
量子计算基础与量子线路 &
量子态、量子门与测量；量子线路模型；普适性与 Clifford+T 综合；量子复杂度基本概念 &
6 & 讲授 + 推导 + 简单线路演示 \\
量子算法基础 &
QFT；相位估计（QPE）；Hadamard test；振幅估计；量子算法的频谱视角 &
9 & 讲授 + 算法推导 + 示例 \\
量子线性代数与后 QPE 方法 &
Block encoding；量子信号处理（QSP）；量子奇异值变换（QSVT）；矩阵函数计算；量子线性系统（HHL） &
9 & 讲授 + 结构分析 \\
量子模拟与动力系统 &
Hamiltonian 模拟；Trotter/LCU 方法；Schr\"odingerization；特征值问题；谱方法与量子模拟的联系 &
9 & 讲授 + 案例分析 \\
微分方程量子算法 &
量子算法求解线性与非线性 PDE（Poisson、波动方程、Hamilton--Jacobi 方程等）；薛定谔化方法；离散化与
矩阵形式；误差来源（离散化/量子误差）；测量复杂性与可观测量计算 &
9 & 讲授 + 文献分析 \\
量子科学计算应用专题（含实践） &
金融（期权定价、风险分析）；材料（特征值问题、电子结构）；机器学习（核方法、回归）；张量网络（tensor networks）；高维问题与量子优势；
UnitaryLab 与 Agent Coding 实践 &
6 & 讲授 + 讨论 + 上机实验 \\
\midrule
\textbf{合计 Total} & & \textbf{48} & \\
\bottomrule
\end{tabularx}

\paragraph{课程思政元素 Ideological \& political elements.}
第 1 章：量子计算作为国家战略产业的意义；第 6 章：量子科学计算的重要应用和价值。

% ==================================================================================
\section*{五、Syllabus (English)}
\addcontentsline{toc}{section}{Syllabus (English)}
\begin{tabularx}{\textwidth}{@{}p{3.2cm} Y c p{3.1cm}@{}}
\toprule
\textbf{Chapter} & \textbf{Main contents} & \textbf{Hrs} & \textbf{Teaching style} \\
\midrule
Fundamentals of Quantum Computing and Quantum Circuits &
Quantum states, gates, and measurement; quantum circuit model; universality and Clifford+T
synthesis; basic concepts of quantum computational complexity. &
6 & Lectures + derivations + simple circuit demos \\
Fundamentals of Quantum Algorithms &
Quantum Fourier Transform (QFT); phase estimation (QPE); Hadamard test; amplitude estimation;
spectral perspective of quantum algorithms. &
9 & Lectures + algorithm derivations + examples \\
Quantum Linear Algebra and Post-QPE Methods &
Block encoding; quantum signal processing (QSP); quantum singular value transformation (QSVT);
matrix function evaluation; quantum linear systems (HHL). &
9 & Lectures + structural analysis \\
Quantum Simulation and Dynamical Systems &
Hamiltonian simulation; Trotter/LCU methods; Schr\"odingerization; eigenvalue problems;
connections between spectral methods and quantum simulation. &
9 & Lectures + case studies \\
Quantum Algorithms for Differential Equations &
Quantum algorithms for linear and nonlinear PDEs (e.g., Poisson, wave, Hamilton--Jacobi
equations); Schr\"odingerization; discretization and matrix formulations; sources of error
(discretization vs quantum); measurement complexity and computation of observables. &
9 & Lectures + literature discussion \\
Applications of Quantum Scientific Computing (with practice) &
Finance (option pricing, risk analysis); materials science (eigenvalue problems, electronic
structure); machine learning (kernel methods, regression); tensor networks; high-dimensional problems and quantum
advantage; UnitaryLab and agent-based coding practice. &
6 & Lectures + discussion + hands-on lab \\
\midrule
\textbf{Total} & & \textbf{48} & \\
\bottomrule
\end{tabularx}

% ==================================================================================
\section*{六、考核方式与要求 \quad Assessment \& Requirements}
\addcontentsline{toc}{section}{考核方式 / Assessment}

\paragraph{中文}
课程考核采用过程性评价与综合评价相结合。平时作业 3--4 次（30\%），考察基础知识与计算能力；
安排 1--2 次小测验（20\%），检验核心概念掌握情况；期末大作业（project，占 50\%），完成一个
综合性项目并提交报告。评分依据正确性、实现质量和表达情况。

\paragraph{English}
The course assessment combines continuous evaluation with a comprehensive final evaluation.
Regular assignments (3--4 in total, 30\%) assess fundamental understanding and computational
skills; 1--2 quizzes (20\%) test mastery of core concepts; and a final project (50\%) requires
completion of a comprehensive task with a written report. Grading is based on correctness,
implementation quality, and clarity of presentation.

\begin{center}
\begin{tabular}{@{}lc@{}}
\toprule
考核项 Component & 占比 Weight \\ \midrule
平时作业 Assignments (3) & 30\% \\
小测验 Quizzes (1) & 20\% \\
期末项目 Final project (report) & 50\% \\ \bottomrule
\end{tabular}
\end{center}

% ==================================================================================
\section*{七、课程资源 \quad Course Resources}
\addcontentsline{toc}{section}{课程资源 / Resources}
\begin{itemize}[leftmargin=1.6em,itemsep=1pt]
  \item 自编讲义 \textit{(self-compiled lecture notes)}.
  \item A.~Childs, \emph{Lecture Notes on Quantum Algorithms}, \url{https://www.cs.umd.edu/~amchilds/qa/}
  \item Lin Lin, \emph{Lecture Notes on Quantum Algorithms for Scientific Computation},
        \url{https://math.berkeley.edu/~linlin/qasc/}
  \item UnitaryLab, \url{https://unitarylab.com/unitary-studio-test/}
\end{itemize}

\end{document}
